This literature review focuses on a narrower question than most retrospective studies of misinformation diffusion: what can actually be learned from the early stage of a cascade, when observation is partial, decision windows are short, and visible diffusion is shaped by attention and platform visibility? Prior work suggests that diffusion structure may contain information beyond early volume, but it is less clear how stable or useful that information remains once the observation window is made explicit. This section therefore reviews work on early diffusion dynamics, mechanisms of spread, and veracity-linked diffusion differences in order to clarify what should reasonably be expected from an early-warning design.



\subsection{Early-Stage Diffusion Dynamics}
\paragraph{Why the early stage is different.}

Prior work on online diffusion and cascade prediction often treats the beginning of a cascade as a distinct analytical regime, because inference must be made under partial observation rather than with complete trajectories \citep{cheng_can_2014, szabo_predicting_2010, colbaugh_earlywarning_2012}. Compared to later stages, early diffusion is usually observed under three constraints: limited information, higher noise, and a short decision window. From a classic diffusion perspective, early adoption is important because diffusion is shaped by social processes and relational structure rather than random exposure, which makes trajectories path-dependent \citep{rogers_diffusion_2003}. In online platforms, this path dependence becomes more visible because the early stage is also where platform ranking, feed curation, and uneven exposure can strongly affect what content becomes observable to a broader audience. A practical reason that early diffusion behaves differently is human attention: users face limited attention and compete across many messages, so what spreads is not only about intrinsic content quality but also about what people actually see and have time to engage with \citep{lerman_information_2010, lerman_information_2016}. This constraint is especially relevant in the early stage: small bursts of engagement can trigger additional exposure, while many items die out quickly because they never enter an attention-amplifying trajectory, so the same ``early'' volume can sometimes mean very different underlying processes.
\paragraph{Defining and evaluating ``early.''}
Given these constraints, prior work has tried to define ``early'' in operational and evaluable terms, typically by setting an explicit observation window and asking what can be predicted using only partial information. In practice, ``early'' is most commonly operationalized in three ways: by elapsed time since the source post, by the first $k$ reshares/replies observed in temporal order, or by partial depth (e.g., the first $d$ generations/levels) of the diffusion tree. In online popularity prediction, early attention measures (such as views or votes) can forecast longer-run popularity, but the useful horizon varies across platforms due to differences in attention decay and user behavior \citep{szabo_predicting_2010}. Work in networked cascade prediction makes the window constraint even more explicit: \citet{cheng_can_2014} formalize an early prediction setup and show that predicting which cascades will become extremely large remains intrinsically difficult, with the heavy-tailed nature of cascades making the tail especially hard to forecast even when more early information is observed. Studies focusing on influence on Twitter further emphasize the same point: \citet{bakshy_everyones_2011} find that highly followed or historically influential users are more likely to generate large cascades, but predicting which specific user or URL will go viral remains unreliable, and this uncertainty is consistent with evidence that network position does not automatically translate into influence when attention is limited \citep{lerman_information_2010}. Taken together, these results suggest a realistic baseline for early-stage inference: early signals can be informative for probabilistic judgments, ranking, or monitoring, but it is risky to expect reliable point prediction of extreme outcomes.


These results clarify what an early-warning objective can realistically be under partial observation. If extreme outcomes are hard to predict as point targets even with more early information \citep{cheng_can_2014}, then the practical goal shifts to decision support under uncertainty: improving relative judgments (ranking, triage, monitoring) within a fixed early window. Framed this way, the key question is whether diffusion structure and shape provide incremental signal beyond simple early activity counts.


\paragraph{Beyond early volume: structure and early warning.}
Under this early-warning objective, an ``early signal'' need not be limited to early volume alone(i.e., simple early activity counts within a fixed observation window). Structural aspects of diffusion can differ even when early counts look similar, and several studies argue that these differences matter for understanding and forecasting. \citet{goel_structural_2016} introduce structural virality to distinguish broadcast-like diffusion from multi-generation spreading, showing that large reach can arise from very different diffusion shapes. Community dispersion provides another structural dimension: \citet{weng_virality_2013} link diffusion outcomes to community structure and show that crossing community boundaries can be predictive of virality, which implies that early breadth and diversity sometimes carry information beyond raw growth rate. Finally, the early-stage constraint matters because interventions are time-sensitive. \citet{colbaugh_earlywarning_2012} explicitly frame diffusion modeling as an early-warning problem where timeliness is part of value, and rumor evidence suggests why: corrective signals can appear after diffusion is already underway, and rumor cascades may continue even in the presence of correction-related activity \citep{friggeri_rumor_2014}. This motivates evaluating early-stage signals under partial observation rather than relying on late-stage information that may be unavailable when decisions are needed.


\subsection{Mechanisms of Spread}

\paragraph{Reinforcement and network structure.}
Mechanism-based work explains how early diffusion differences can arise from reinforcement (complex contagion) and community structure, rather than from aggregate growth alone. Classic diffusion theory emphasizes that adoption is socially structured: individuals adopt ideas and behaviors through communication channels and social systems, not in isolation \citep{rogers_diffusion_2003}. One influential mechanism is reinforcement. In an online network experiment, \citet{centola_spread_2010} show that behaviors requiring social reinforcement spread more effectively in clustered networks, capturing the logic of complex contagion where multiple exposures are often necessary for adoption. In online diffusion settings, this mechanism provides a reason to treat local structure and repeated exposure as meaningful: early clustering and density can either stabilize diffusion through reinforcement or constrain it within a small neighborhood. Community boundaries play a similar mechanistic role. Empirical evidence suggests that diffusion outcomes depend on whether a cascade stays inside a community or crosses communities, and \citet{weng_virality_2013} show that early dispersion across communities can be associated with larger cascades. Together, reinforcement and community structure offer a relatively concrete account of why early diffusion patterns may differ across cascades even before considering content.

\paragraph{Attention constraints and diffusion shape.}
At the same time, online diffusion is constrained by attention and cognition in ways that complicate simple epidemic metaphors. \citet{lerman_information_2016} argues that ``information is a virus'' can be misleading because exposure is limited, users face intense competition among messages, and platform filtering shapes what people see. Earlier empirical work on Digg and Twitter supports this view, showing that most content does not spread far and that influence is not purely a function of centrality (how well-connected or structurally important a user is in the network) \citep{lerman_information_2010}. These constraints interact with diffusion shape. Some large cascades are shallow and broadcast-driven while others are deeper and multi-step, and \citet{goel_structural_2016} show that these forms coexist at scale; influence-focused evidence similarly suggests that who shares matters without making outcomes fully predictable \citep{bakshy_everyones_2011}. Mechanistically, this implies that early diffusion signals can reflect a mix of social reinforcement, cross-community reach, and visibility/attention shocks, which is one reason early-stage inference is both useful and noisy.

\paragraph{Polarization, selective exposure, and veracity.}
Mechanisms become more specific in the misinformation literature, where polarization, selective exposure, and uneven participation are central. \citet{del_vicario_spreading_2016} document echo chambers and selective exposure on Facebook, showing that misinformation diffusion is shaped by homophily and polarization, which can keep diffusion concentrated within like-minded groups. On Twitter during the 2016 U.S. election, \citet{grinberg_fake_2019} found that fake news exposure and sharing are highly concentrated among a small subset of users, suggesting strong heterogeneity in who participates in misinformation diffusion. Large-scale evidence also points to systematic differences by veracity: \citet{vosoughi_spread_2018} show that false news spreads faster, farther, and deeper than true news on Twitter and argue these differences are not primarily driven by bots. Broader interdisciplinary and economic framings emphasize that misinformation is not a single-mechanism story; incentives, psychology, networks, and platform design all matter \citep{allcott_gentzkow_2017, lazer_science_2018}. These perspectives help interpret diffusion patterns as behavioral and contextual outcomes rather than purely content-intrinsic properties of content. These mechanisms matter for how early signals should be interpreted. If reinforcement, community boundaries, and attention constraints shape who gets exposed and who participates, then early diffusion is not only about eventual reach: the same early growth can reflect different social environments, and those environments may plausibly differ between reliable and misleading content, especially in polarized or highly clustered settings. Taken together, these mechanisms suggest that early diffusion is shaped by reinforcement, community structure, and attention constraints; the next section asks whether these forces produce veracity-linked diffusion differences that emerge early enough to be useful under explicit early-observation windows.

\subsection{Veracity Signals and Implications for Early Warning}


\paragraph{Early predictability of reach.}
The diffusion prediction and misinformation literatures are often linked by a simple idea: if misinformation spreads differently, then early diffusion patterns might be useful signals \citep{shu_fake_2017,zubiaga_detection_2018,vosoughi_spread_2018}. On the diffusion side, a recurring theme is that early information helps more with coarse judgment (ranking/monitoring) than with precise point prediction of rare, extremely large cascades. Early attention can forecast longer-run popularity, but how long an ``early'' window remains informative depends on platform dynamics and attention decay \citep{szabo_predicting_2010}. Formal early-prediction setups reach a similar conclusion: under heavy-tailed outcomes, the tail is hard to forecast even as more early information is observed \citep{cheng_can_2014}. Influence studies also point to limits of specificity: influential users are more likely to generate large cascades on average, but which particular user or URL will go viral remains hard to predict \citep{bakshy_everyones_2011}, and network position does not translate cleanly into influence when attention is limited \citep{lerman_information_2010}. This is where structural descriptors enter: different diffusion shapes (shallow broadcast-like vs.\ deeper multi-step spreading) can imply very different reach profiles \citep{goel_structural_2016}, and early dispersion across communities can add information beyond raw volume \citep{weng_virality_2013}. This motivates evaluating early structural and temporal features under explicit early-observation windows.

 
\paragraph{Retrospective evidence on veracity-linked diffusion.}
On the misinformation side, there is strong descriptive evidence that diffusion differs by veracity, but much of it is retrospective. \citet{vosoughi_spread_2018} compare complete cascades and document that false news spreads faster and farther, while \citet{del_vicario_spreading_2016} document echo chambers and polarization that shape misinformation diffusion on Facebook. In the election context, \citet{grinberg_fake_2019} show that exposure and sharing can be highly concentrated among a small subset of users, reinforcing the idea that diffusion reflects both network structure and participation asymmetry. Rumor cascade evidence illustrates the timing problem: correction-related signals can appear after diffusion is already underway, and rumors can continue spreading even with correction-related activity \citep{friggeri_rumor_2014}. This is why early-warning framing matters. \citet{colbaugh_earlywarning_2012} emphasize that timeliness is part of value, which suggests that methods relying on late-stage information may be accurate but not actionable.

\paragraph{Detection methods and early-window evaluation.}
Building on the mechanism-based perspective above, this section turns to how misinformation detection studies operationalize and evaluate \emph{early} settings, since evaluation design largely determines what claims about early warning can be supported. Detection-focused research provides a map of methods but also reveals evaluation challenges in early-stage settings. Survey work organizes rumor and fake news research into pipelines and feature categories and highlights early-stage uncertainty and limited ground truth \citep{shu_fake_2017, zubiaga_detection_2018}. Model-based studies show that temporal and propagation signals can help detection: \citet{ma_detecting_2016} model temporal dynamics in microblogs, while propagation-aware approaches leverage structure more directly, including tree-structured recursive models \citep{ma_tree_rnn_2018} and graph convolutional networks with directional information flow \citep{bian_bigcn_2020}. These results support the claim that propagation structure contains veracity-relevant information, but they also make clear why comparability across studies is difficult: many evaluations assume substantial portions of a cascade are observable, and ``early'' is sometimes defined loosely. The early-prediction literature provides a more disciplined window-based evaluation approach \citep{cheng_can_2014}, while early-warning work emphasizes timeliness and partial observability \citep{colbaugh_earlywarning_2012}. Connecting these norms to misinformation detection remains an important methodological issue, especially given that diffusion signals reflect multiple interacting forces---reinforcement, community boundaries, polarization, and attention constraints \citep{centola_spread_2010, lerman_information_2016, del_vicario_spreading_2016}---and are therefore context-dependent rather than purely intrinsic properties of content. Against this background, the literature leaves a fairly concrete early-warning question. On the one hand, early diffusion signals can support partial forecasting of eventual reach, but prediction is intrinsically limited and depends on how the early window is defined \citep{szabo_predicting_2010, cheng_can_2014, bakshy_everyones_2011}. On the other hand, veracity is associated with systematic differences in diffusion outcomes when cascades are viewed retrospectively \citep{vosoughi_spread_2018, del_vicario_spreading_2016}. What is less clearly established is when such veracity-linked differences become detectable under realistic early windows defined by time, volume, or depth, and how this visibility trades off with predictive performance. Taken together, this motivates the design choices in the next section: I define explicit early-observation windows and test whether diffusion structure adds value beyond simple early activity for reach and veracity under partial observation.

Accordingly, this study asks whether early diffusion patterns---structural and temporal signals extracted from the initial portion of a cascade---can predict (i) eventual spread/reach and (ii) veracity labels in tree-structured rumor datasets. The next sections specify the datasets, the early-window definitions, and the feature construction used to test these questions under partial observation.

\subsection{Hypotheses}
The literature above leads to two directional expectations. First, early diffusion signals should be more informative for eventual reach than for veracity, because forecasting later spread is more directly aligned with early cascade growth than judging truth status from partial propagation alone.

\textbf{H1.} Early diffusion features will predict eventual reach better than they predict veracity under the same early-observation window.

Second, if structural shape contains information beyond early activity alone, then richer structural and dynamic bundles should improve performance over tempo-only baselines, especially for the reach task.

\textbf{H2.} Structural and dynamic diffusion features will add more predictive value for reach than for veracity beyond simple early-activity baselines.
